\section{Introduction}\label{Introduction}

In \cite{FG1}, Fock and Goncharov define a class of schemes, called cluster varieties, whose rings of global regular functions are upper cluster algebras. In \cite{GHK3}, Gross, Hacking, and Keel describe how to view cluster varieties as certain blowups of toric varieties. We review this description, as well as Gross--Hacking--Keel's construction \cite{GHK1} of the tropicalization of a log Calabi--Yau surface. We then use these ideas to give a classification of rank\footnote{Cluster algebraists often take rank $2$ to mean that the exchange matrix is $2\times 2$. However, we use rank to mean the dimension of the space spanned by the rows or columns of the exchange matrix.} $2$ cluster varieties (those for which the symplectic leaves of the $\s{X}$-space are $2$ dimensional) and to describe their cluster modular groups. This can also be viewed as a classification of log Calabi--Yau surfaces.

By a {\it log Calabi--Yau surface} or a {\it Looijenga interior}, we mean a surface $U$ which can be realized as $Y\setminus D$, where $Y$ is a smooth, projective, rational surface over an algebraically closed field $\kk$ of characteristic~$0$, and the {\it boundary}~$D$ is a choice of snc anti-canonical divisor in $Y$. Furthermore, $D=D_1+\cdots + D_n$ is either a cycle of smooth irreducible rational curves $D_i$ with normal crossings, or if $n=1$, $D$~is an irreducible curve with one node. By a {\it compactification} of~$U$, we mean such a pair $(Y,D)$ (\cite{GHK2} calls these compactifications with ``maximal boundary''). We call $(Y,D)$ a {\it Looijenga pair}, as in~\cite{GHK1}. Toric varieties are the most basic examples, and every~$U$ can be obtained by performing certain blowups on a toric surface, cf.\ Lemma~\ref{ToricModel}.

